\documentclass[11pt,a4paper]{article}
\usepackage[margin=27mm]{geometry}
\usepackage{amsmath,amssymb,amsthm,fontspec,unicode-math,hyperref,booktabs,longtable}
\usepackage{array,calc,fvextra}
\setmainfont{STIXGeneral}
\setmathfont{STIX Two Math}
\setmonofont{Menlo}
\hypersetup{hidelinks}
\setlength{\emergencystretch}{2em}
\providecommand{\tightlist}{\setlength{\itemsep}{0pt}\setlength{\parskip}{0pt}}
\newtheorem{theorem}{Theorem}
\newtheorem{lemma}[theorem]{Lemma}
\newtheorem{corollary}[theorem]{Corollary}
\newcommand{\gen}[1]{\langle #1\rangle}
\newcommand{\Sym}{\operatorname{Sym}}
\newcommand{\Aut}{\operatorname{Aut}}
\newcommand{\Fix}{\operatorname{Fix}}

\title{Regular polyhedra with an odd kernel of index four}
\author{The human\_compression project}
\date{Canonical Topic F: application}
\begin{document}
\maketitle
\begin{abstract}
For every odd prime $p$, we recover the structural and abelian-kernel
classification of rank-three string C-groups of order $4p^m$ and give
their nonabelian family in cyclotomic coordinates. The construction is
identified generator by generator with the original presentation.
The argument separates a two-character classification from a ring
calculation valid for all parameters. A literal marking in the specified
primary edition is corrected explicitly.
\end{abstract}
\section{The structural theorem}
We use the twelve-page preprint arXiv:2407.10388v1 by Hou, Feng,
Leemans and Qu~\cite{HFLQ}. We recover its Theorems 1.1--1.3
and the explicit construction in Theorem 5.1, with the correction stated
below. The results remain attributed to those authors. This edition
reorganizes the accepted proofs through the common paper~\cite{UF};
no previous combined manuscript is an additional proof premise.

We use the conventions of that paper: marked generators are distinct
nonidentity involutions, the endpoints commute, and the adjacent
orders are at least three. A string C-group satisfies the full subset
intersection property. Its type is $\{k_1,k_2\}$, and it is tight
when $|G|=2k_1k_2$. Conjugation is $x^g=g^{-1}xg$.

% BEGIN F1.Structure
\begin{theorem}[The odd kernel]\label{F1.Structure}
Let $p$ be an odd prime and $m\geq2$. A nondegenerate rank-three
string C-group $(G;r_0,r_1,r_2)$ of order $4p^m$ has a normal
Sylow $p$-subgroup $P$, and
\[
 G=P\rtimes K,\qquad K=\gen{r_0,r_2}\cong C_2\times C_2.
\]
Both type parameters are divisible by $p$, at least one is divisible
by $2p$, and $P$ needs exactly two generators.
\end{theorem}
\begin{proof}
The number $n_p$ of Sylow $p$-subgroups divides four and is congruent
to one modulo $p$. The only case other than $n_p=1$ is
$p=3$, $n_p=4$. Assume this case occurs. Conjugation on the four
Sylow subgroups gives a homomorphism $G\longrightarrow S_4$ with
kernel $N\leq N_G(P)=P$, since the normalizer has index four.
The image has order $4\cdot3^j$ dividing $24$, so $j=0$ or $1$.
If $j=0$, then $N=P$ and $P$ is normal, a contradiction. If $j=1$,
the image is the index-two subgroup $A_4$ of $S_4$. Its involutions
all lie in the normal Klein four subgroup. They cannot generate
$A_4$, whereas an image of an involution-generated group is generated
by involutions and identities. This is also a contradiction.
Thus $P$ is normal.

The endpoint subgroup $K$ has order four, since its two generators
are distinct commuting involutions. It intersects the odd group $P$
trivially, and $|PK|=|G|$, giving the stated semidirect product.
Each adjacent dihedral group has order $2k_i$, so
$k_i\mid2p^m$. Nondegeneracy implies $p\mid k_i$.
If both $k_i$ were odd, their images in the exponent-two quotient
$G/P=K$ would force $r_0P=r_1P=r_2P$, contrary to the distinct
images of $r_0,r_2$. Thus one label is even.

The inverted-orbit lemma of~\cite{UF} gives two generators for $P$.
If $P$ were cyclic, the nontrivial $p$-parts of both adjacent rotation
groups would contain its unique subgroup of order $p$. That subgroup
would lie in $G_{\{0,1\}}\cap G_{\{1,2\}}=\gen{r_1}$,
which has order two. This contradiction proves that two generators
are necessary.
\end{proof}
% END F1.Structure

\section{Abelian kernels and tight markings}
% BEGIN F1.Abelian
\begin{theorem}[The abelian classification]\label{F1.Abelian}
Under the hypotheses of Theorem~\ref{F1.Structure}, the kernel $P$
is abelian if and only if the marking is tight. Up to duality, all
such markings are the following, one for each $1\leq a\leq b$,
$a+b=m$:
\begin{gather*}
 P=\gen{x}\times\gen{y},\qquad |x|=p^b,\quad |y|=p^a,\\
 x^{r_0}=x,\quad y^{r_0}=y^{-1},\qquad
 x^{r_2}=x^{-1},\quad y^{r_2}=y^{-1},\qquad r_1=xy r_2.
\end{gather*}
The type is $\{2p^a,p^b\}$. For this type the marked group is unique.
\end{theorem}
\begin{proof}
First suppose $P$ is abelian. Write $r_1=ut$ with $u\in P$ and
$1\ne t\in K$. The two-character lemma of~\cite{UF} shows that
$P$ is the direct product of the two nonzero cyclic character
components on which $t$ acts as $-1$. Choose their projected
generators $x,y$ so that $u=xy$.

If $t=r_0r_2$, the two characters on $(r_0,r_2)$ are $(+,-)$ and
$(-,+)$. A direct product calculation gives adjacent orders
$2|y|$ and $2|x|$. The rank-three product bound of~\cite{UF} would
give $8|x||y|\leq |G|=4|x||y|$, which is impossible.
Thus $t$ is an endpoint; by duality take $t=r_2$. The two characters
are now $(+,-)$ and $(-,-)$, exactly as in the displayed model.
For $|x|=p^b$ and $|y|=p^a$ the adjacent orders are
\[
 |r_0r_1|=2p^a,\qquad |r_1r_2|=p^{\max(a,b)}.
\]
The same bound implies $a+\max(a,b)\leq a+b$, hence $a\leq b$.
It is then an equality; the marking is tight.

Conversely, construct the displayed semidirect product for
$1\leq a\leq b$. The two specified automorphisms generate a
Klein four complement, and $r_1$ is an involution because $r_2$
inverts $xy$. The three marked elements are distinct, the endpoints
commute, and the adjacent orders are $2p^a,p^b$.
The elements $u=xy$ and $u^{r_0}=xy^{-1}$ yield $x^2,y^2$;
squaring is invertible on both cyclic odd-order factors. Hence the
marking generates $G$.

Put $H=\gen{r_0,r_1}$ and $J=\gen{r_1,r_2}$.
The image of $H$ in $K$ has order four. Its rotation square
generates $\gen{y}$, and $|H|=4p^a$, so $H\cap P=\gen{y}$.
Similarly $J\cap P=\gen{xy}$ and its image in $K$ is
$\gen{r_2}$. These two odd subgroups meet trivially: if $x^j=1$
then $p^b\mid j$, which also gives $y^j=1$ since $a\leq b$.
The map $H\cap J\longrightarrow K$ is consequently injective,
with image in $\gen{r_2}$. It contains $r_1$, so
$H\cap J=\gen{r_1}$. The rank-three test proves the full
intersection property. The cyclic factor orders, their projected
generators and the two characters determine the marked group,
giving uniqueness among abelian-kernel models.

It remains to show that tightness implies an abelian kernel. By
Theorem~\ref{F1.Structure}, each type parameter is a power of $p$
or twice a power of $p$. The equality $|G|=2k_1k_2=4p^m$
forces exactly one parameter to be even, and the exponents sum to
$m$. We now retain the precise odd/even clause of
Cunningham--Pellicer's classification~\cite[Theorem 1.1(b),(c)]{CP}:
apart from its exceptional even parameter four, a tight regular
polyhedron with one odd and one even type parameter exists only
when the even parameter divides twice the odd parameter, and the
marked polyhedron is unique for such a type. Here the even parameter
is $2p^a$, never four. Thus $a\leq b$, and uniqueness identifies
the marking with the abelian model just constructed. The use of that
theorem includes the retained regular-polyhedron/string-C-group
correspondence. This completes the proof.
\end{proof}
% END F1.Abelian

% BEGIN F1.Correction
\paragraph{The literal formula in the primary edition.}
The display in~\cite[Theorem 1.2, physical p.\ 2]{HFLQ} gives
$x^{r_0}=x$ and $y^{r_0}=y^{-1}$ but writes $r_1=xy r_0$.
Since $x,y$ commute, this element has square $x^2\ne1$.
It is not an involution for the asserted positive odd-prime exponent.
The body of the same edition, Section 4 on physical page 6, uses
$xy r_2$ in the corresponding character case. Theorem~\ref{F1.Abelian}
uses that formula and assigns the larger exponent to $x$.
This is a correction to the specified preprint display, not a recovery
of its impossible literal formula or an assertion about a later
journal edition.
% END F1.Correction

\section{A nonabelian kernel for every parameter}
% BEGIN F1.Construction
\begin{theorem}[The cyclotomic polyhedra]\label{F1.Construction}
For every odd prime $p$ and every $m\geq3$, there is a string
C-group of order $4p^m$ and type $\{p,2p\}$ with nonabelian
Sylow $p$-subgroup. More precisely, let
\[
 O=\mathbb Z[\zeta_p],\quad \pi=\zeta_p-1,\quad
 A=O/\pi^{m-1}O,\quad P=A\rtimes\gen{\beta},
 \qquad |\beta|=p,\quad v^\beta=\zeta_pv.
\]
Define automorphisms by
\[
 \sigma(v)=\overline v,\quad\sigma(\beta)=\beta^{-1},
 \qquad\tau(v)=-v,\quad\tau(\beta)=\beta.
\]
Then $G=P\rtimes\gen{\sigma,\tau}$ has the required marking
\[
 r_0=s\tau\sigma,\qquad r_1=\beta\tau\sigma,\qquad r_2=\sigma,
\]
where $s$ denotes the additive class of the ring element $1$.
It is not the identity of the additive group $A$.
\end{theorem}
\begin{proof}
This is the shared split cyclotomic construction in the accompanying
common paper~\cite{UF}, with $n=m-1$. That paper includes the complete
normal-form, automorphism, generation and intersection proof.
\end{proof}
% END F1.Construction

\section{The exact presentation and its automorphisms}
% BEGIN F1.Presentation
\begin{theorem}[The source-generator identification]\label{F1.Presentation}
Write
\[
 m-1=(e-1)(p-1)+r,\qquad e\geq1,\quad1\leq r\leq p-1,
\]
and put $d_i=p^e$ for $i\leq r$ and $d_i=p^{e-1}$ for $i>r$.
The Sylow subgroup in Theorem~\ref{F1.Construction} has the exact
presentation
\begin{align*}
 \langle s_1,\ldots,s_{p-1},\beta\mid{}
 &s_i^{d_i}=1,\quad\beta^p=1,\quad [s_i,s_j]=1,\\
 &[s_i,\beta]=s_{i+1}\ (1\leq i<p-1),\\
 &[s_{p-1},\beta]=\prod_{i=1}^{p-1}s_i^{-\binom pi}\rangle.
\end{align*}
The power and commutativity relators range over $1\leq i,j\leq p-1$.
The generator $s_i$ maps to $\pi^{i-1}$ in $A$ and has exactly
order $d_i$. Under this isomorphism, $\sigma$ fixes $s_1$ and
inverts $\beta$, while $\tau$ inverts $s_1$ and fixes $\beta$.
Thus the group, automorphisms and marking are precisely those of
\cite[Section 5 and Theorem 5.1]{HFLQ}.
\end{theorem}
\begin{proof}
The coordinate theorem of~\cite{UF} gives the exact additive orders
and commutativity of the proposed $s_i$. Our commutator convention
gives $[s_i,\beta]=\pi s_i=s_{i+1}$ for $i<p-1$.
The cyclotomic equation
\[
 p+\binom p2\pi+\cdots+\binom p{p-1}\pi^{p-2}
       +\pi^{p-1}=0
\]
gives the final relation. Thus all literal source relators hold.

Conversely let $\widehat P$ be the group defined by the displayed
presentation. Its subgroup $S=\gen{s_1,\ldots,s_{p-1}}$ is finite
abelian, with order at most
$\prod_i d_i=p^{er+(e-1)(p-1-r)}=p^{m-1}$.
The conjugation relators give $S^\beta\leq S$; since conjugation
is injective and $S$ is finite, this is equality. Hence $S$ is normal
and $\widehat P/S$ is cyclic of order at most $p$.
Consequently $|\widehat P|\leq p^m$. The homomorphism to the
constructed group $P$ is surjective: the displayed basis generates
$A$, and $\beta$ generates the quotient. Its image has order
$p^m$, so the homomorphism is an isomorphism. This proves the
generator-by-generator assertion, including modulus-one generators.

The formulas on $s_1,\beta$ are the definitions of $\sigma,\tau$.
These elements generate $P$, because successive commutators of
$s_1$ with $\beta$ give every $s_i$. The formulas therefore identify
the automorphisms, not only their effects on a proper subgroup.
\end{proof}
% END F1.Presentation

% BEGIN F1.Binomial
\paragraph{The auxiliary binomial formulas.}
Define $s_p=[s_{p-1},\beta]$ and $s_{p+1}=[s_p,\beta]$.
They are the additive classes $\pi^{p-1},\pi^p$; they are not
additional independent generators. The common conjugation identities
give, in multiplicative notation, for $1\leq k\leq p$,
\begin{align*}
 s_k^{\beta^{1-k}}&=\prod_{i=k}^{p+1}
           s_i^{\binom{p+1-k}{i-k}},\\
 s_k^\sigma&=
 \begin{cases}
 s_k^{\beta^{1-k}},&k\text{ odd},\\
 (s_k^{\beta^{1-k}})^{-1},&k\text{ even}.
 \end{cases}
\end{align*}
These are exactly Lemma 5.2(1),(2) of~\cite{HFLQ}.
The verified automorphisms in Theorem~\ref{F1.Construction} give
Lemma 5.3 directly. The source's coefficient-table calculation is
therefore unnecessary for this proof: the full relations are
preserved by actual automorphisms of the ring model and the exact
presentation isomorphism transports them back.
% END F1.Binomial

% BEGIN F1.Class
\paragraph{Nilpotency class and the Hughes hypothesis.}
The group $A$ is abelian, $P/A$ is cyclic, and
$[A,\beta]=\pi A$. Induction gives
\[
 \gamma_j(P)=\pi^{j-1}A\qquad(j\geq2).
\]
The successive lattice indices show that these groups are nontrivial
through $j=m-1$ and trivial at $j=m$. Thus $P$ has nilpotency
class $m-1$, the maximal-class property stated in the source.
The motivating condition $A=H_p(P)$ in its Section 5 is not a
condition needed for the constructed presentation. In fact, if
$m-1\leq p-1$, the coordinate theorem makes $A$ elementary
abelian. All elements outside $A$ have order $p$ as well, so every
element of $P$ has order dividing $p$. The Hughes subgroup, generated
by elements of order different from $p$, is then trivial. No general
classification of Hughes groups or groups of maximal class is used
in any of our arguments.
% END F1.Class

\section{Source scope and verification}
% BEGIN F1.Scope
Theorem~\ref{F1.Structure} is the structural result
\cite[Theorem 1.1]{HFLQ}. Theorem~\ref{F1.Abelian} recovers its
Theorem 1.2 with the explicit marking correction. Theorem~\ref{F1.Construction}
and the exact presentation recover Theorems 1.3 and 5.1. The auxiliary
rank-three criterion, character decomposition, binomial identities,
automorphisms and maximal-class assertion are included. The original
small-group numerical census, survey discussion and general Hughes
classification are outside this principal-results contract.

The mathematical results retain their original attribution. The common
inverted-orbit and two-character arguments reorganize the source proofs;
the cyclotomic module supplies a uniform replacement for the construction
and its presentation checks. We make no novelty claim for the underlying
cyclotomic model. In the tightness converse we retain exactly
\cite[Theorem 1.1(b),(c)]{CP} in the stated odd/even range. That
primary theorem and the polyhedron correspondence were read in the
earlier accepted audit, but its full PDF is not retained locally.
There is no local byte binding or known complete-text cost for that
external proof. It is not silently replaced by finite computations.

The complete accompanying validation appendix retains three programs:
the mixed-modulus model generator, the finite marked-group checker,
and the independently authored literal-presentation checker. Together
they test seven cyclotomic and four abelian marked models and five
literal source presentations. Every tested marking is checked against
all 64 subset-intersection identities. The programs also reject the
incorrect literal marking, a corrupted cyclotomic relation and altered
source generator orders; a composite-index example shows that the
prime-index hypothesis is necessary. These are bounded checks of the
uniform proofs above, whose conclusions do not depend on a finite
enumeration or an unreported program output. The three complete
programs remain visible and chargeable as validation; both complete
historical operational drivers are inventoried separately.

The five reused Lean declarations verify the two prime-index interfaces
and three group identities. Sylow normality, the two-character theorem,
cyclotomic lattice calculations, all-parameter construction and
source-presentation isomorphism remain written proofs. The retained
tight-polyhedron classification and classical correspondence are also
outside that formal scope. This edition neither solves a new open
question nor changes the previously admitted source count. It asserts
no compression ratio; the common paper, complete application, validation
programs and external inputs have separate, explicit boundaries.
% END F1.Scope

\begin{thebibliography}{9}
\bibitem{UF} The human\_compression project,
\emph{Odd kernels, cyclotomic modules and cyclic carries}, accompanying common paper.
\bibitem{HFLQ} D.-D. Hou, Y.-Q. Feng, D. Leemans and H.-P. Qu,
\emph{String C-groups of order $4p^m$}, arXiv:2407.10388v1,
15 July 2024, 12 pp.
\bibitem{CP} G. Cunningham and D. Pellicer,
\emph{Classification of tight regular polyhedra}, arXiv:1408.6703v2,
10 November 2015, Theorem 1.1(b),(c); published in
\emph{Journal of Algebraic Combinatorics} 43 (2016), 665--691.
\end{thebibliography}
\end{document}
